\chapter{Path Deduplication and Pagination}
\label{ch:dedup-pagination}

The matcher can return many paths for a lax hypothesis. Before an
analyst sees them they're deduplicated (depending on the requested
mode), scored, filtered, sorted, and paginated.
\code{score\_and\_paginate\_paths} is the single function that does
all of that; it lives in \filepath{graph\_hunter\_core/src/analytics.rs}.

\section{Dedup modes}

\begin{definition}[Dedup modes]
\code{DedupMode} is one of:
\begin{itemize}
    \item \code{None} --- every path from the matcher is kept.
    \item \code{ByPath} --- paths with identical vertex sequences
        are merged into a single representative (the highest-scoring
        one).
    \item \code{ByEndpoints} --- paths sharing the same $(v_0,
        v_\ell)$ pair are merged, regardless of the intermediate
        vertices.
\end{itemize}
\end{definition}

\code{ByPath} is the default for hunts rendered in the UI: when the
underlying graph has many parallel edges between the same pair of
vertices, the matcher produces multiple identical-looking paths
(one per edge), and the analyst only wants to see the topology
once.

\code{ByEndpoints} is used during bulk analysis: ``how many
distinct start-end reachability pairs satisfy this hypothesis,''
which is the right granularity for counting unique attack
scenarios.

\section{Algorithm}

\begin{algorithm}
\caption{\textsc{Score-And-Paginate}}\label{alg:score-paginate}
\KwIn{paths $\Pi$, page index $p$, page size $s$, optional
      $\text{min\_score}$, mode $M$, scorer $S$}
\KwOut{page of paths + total count (pre-page)}
\tcp{Step 1: dedup according to mode}
\Switch{$M$}{
    \Case{\code{None}}{unique $\gets \Pi$}
    \Case{\code{ByPath}}{
        key each $\pi \in \Pi$ by its vertex tuple; keep the
        highest-scoring $\pi$ per key
    }
    \Case{\code{ByEndpoints}}{
        key each $\pi \in \Pi$ by $(v_0, v_\ell)$; keep the
        highest-scoring $\pi$ per key
    }
}
\tcp{Step 2: score}
\ForEach{$\pi \in $ unique}{
    $(\text{composite}, \text{breakdown}) \gets S.\text{score\_path}(\pi)$
}
\tcp{Step 3: optional filter}
\If{$\text{min\_score}$ is Some}{drop $\pi$ with composite $<$ min\_score}
\tcp{Step 4: sort DESC by (composite, max\_node\_score, total\_score)}
sort unique\;
\tcp{Step 5: slice}
$\text{total} \gets |\text{unique}|$\;
\Return (unique[$p \cdot s$ .. $(p+1) \cdot s$], total)\;
\end{algorithm}

\begin{complexity}
Let $P = |\Pi|$ (raw count) and $U = |\text{unique}|$ after dedup.
\begin{itemize}
    \item Dedup in \code{ByPath}/\code{ByEndpoints} is
        $\Ord{P \cdot \ell}$ (hashing the key of length $\ell$);
        \code{None} is $\Ord{P}$.
    \item Scoring is $\Ord{U \cdot \ell}$ using \textsc{Score-Path}.
    \item Sorting is $\Ord{U \log U}$ comparisons.
    \item Slicing is $\Ord{s}$.
\end{itemize}
Total: $\Ord{P \cdot \ell + U \log U + s}$. Space is $\Ord{U + s}$
above the input.
\end{complexity}

\section{Sort stability}

The sort key is a tuple
\[
    (\mathrm{composite}(\pi), \max_{v \in \pi}\mathrm{composite}(\{v\}),
     \mathrm{total\_score}(\pi))
\]
sorted in descending order. The three-way key matters because
composite scores tie frequently on paths with short length ---
secondary and tertiary keys break ties deterministically. In
particular the max-node score promotes paths that contain a
single strongly anomalous vertex over paths that are uniformly
mediocre.

\section{Pagination contract}

\begin{invariant}[Page stability]
Two consecutive \code{get\_hunt\_page} calls with the same
parameters return the same sorted order, provided the underlying
graph has not been mutated. Path order is deterministic because
every tie-breaker is a totally ordered quantity derived from the
graph state.
\end{invariant}

The parity harness (Chapter~\ref{ch:parity-harness}) pins this
invariant with a fixture that runs a hunt twice and asserts
element-wise equality of the returned page.

\begin{inthecode}
\filepath{graph\_hunter\_core/src/analytics.rs:1096} ---
\code{score\_and\_paginate\_paths}.\\
\filepath{graph\_hunter\_core/src/analytics.rs} --- \code{DedupMode}
variants and key derivation helpers.
\end{inthecode}
